\pagebreak
\section{Исходный код}
Программа \texttt{main.py} состоит из регулярного шаблона \texttt{PATTERN},
функции \texttt{tokenize} (токенизация и нормализация), функции \texttt{baseline}
(наивное разбиение для сравнения) и функции \texttt{run} (прогон по корпусу и
статистика). Приведён упрощённый шаблон (кириллические классы символов заменены
на латинские для компактности; адрес почты опущен) и основная функция.

\begin{lstlisting}[language=Python]
PATTERN = re.compile(r"""
    (?P<url>(?:https?://|www\.)[^\s<>"'\])]*[^\s<>"'\]).,;:!?]) |
    (?P<initials>[A-HJ-UW-Z]\.\s?(?:[A-HJ-UW-Z]\.\s?)?(?=[A-Z][a-z])) |
    (?P<abbr>
        (?i:(?:t\.\s?[dpe]|n\.\s?e)\.) |          # "t. d.", "n. e."
        (?:[^\W\d_]\.){2,} |                       # dotted: "t.d.", "e.g."
        (?:(?<=[\dIVX][ \xa0])|(?<=[\dIVX]))[vg]\. |   # "XX v.", "1990 g."
        (?i:(?:%s)\.)                              # dictionary of abbreviations
    ) |
    (?P<number>
        (?<!\d)\d{1,3}(?:[ \xa0\u202f\u2009]\d{3}(?!\d))+(?![\w-]) |  # 10 000
        \d+(?:[.,:/]\d+)*(?!\w|-\w)                                   # 3,14  12.03.2020
    ) |
    (?P<compound>[^\W_]+(?:[-\u2010\u2011][^\W_]+)+) |  # word-word, F-16, 1990-e
    (?P<word>[^\W_]+(?:['\u2019][^\W_]+)*) |            # words, O'Neil, 3D
    (?P<punct>[^\w\s])
""" % "|".join(ABBR), re.VERBOSE)

DROP = str.maketrans("", "", "\u00ad\u0301")  # soft hyphen, stress mark

def tokenize(text, lower=True, yo=True, punct=False, split_compounds=False):
    out = []
    for m in PATTERN.finditer(text.translate(DROP)):
        kind, tok = m.lastgroup, m.group()
        if kind == "punct" and not punct:
            continue
        if kind in ("abbr", "initials"):
            tok = re.sub(r"\s", "", tok)          # "t. e." -> "t.e."
        if kind != "url" and kind != "email":
            tok = tok.lower() if lower else tok
            tok = tok.replace("\u0451", "e") if yo else tok
        if kind == "compound" and (split_compounds
                or re.fullmatch(r"[\d\-\u2010\u2011]+", tok)):
            # numeric ranges are never one token
            out += [("number" if p.isdigit() else "word", p)
                    for p in re.split("[-\u2010\u2011]", tok)]
        else:
            out.append((kind, tok))
    return out
\end{lstlisting}

В реальном коде классы символов и словарь сокращений содержат кириллицу
(\texttt{А-ЯЁ}, \texttt{т}, \texttt{д} и т.\,д.), а \enquote{ё} записана
буквой; в листинге это заменено на латиницу и код \texttt{\textbackslash u0451}, так как
\texttt{listings} плохо работает с кириллицей.

\pagebreak
\section{Выводы}
Выполнив третью лабораторную работу по курсу \enquote{Информационный поиск}, я
написал токенизатор для русскоязычных текстов и проверил его на корпусе из
30\,000 статей Википедии. Основное, что я узнал нового:

\begin{itemize}
    \item токенизация --- не \enquote{разбить по пробелам}: наивное разбиение по
    не-буквам даёт 18{,}29 млн токенов, а токенизатор с правилами --- 17{,}77 млн.
    Разница в 0{,}5 млн токенов --- это числа, даты, сокращения, инициалы и
    составные слова, которые наивный метод рвёт на мусор из одиночных букв и цифр;
    \item порядок правил важен: регулярное выражение берёт первый подходящий
    вариант, поэтому ссылка, инициалы, сокращение и число должны стоять раньше
    слова, а число --- проверять, что за ним не идёт буква или дефис, иначе
    \enquote{3D} и \enquote{1990-е} распадаются;
    \item неоднозначность решается эвристиками: точка после \enquote{г} --- сокращение,
    но \enquote{в} --- предлог, поэтому \enquote{в.} признаётся сокращением только после
    числа или римской цифры; \enquote{Пётр I. Он} нельзя принять за инициалы, поэтому
    римские цифры исключены; десятичную запятую (\texttt{3,14}) и разделитель тысяч
    (\texttt{10 000}) нужно отличать от списка чисел, поэтому группа обязана
    состоять ровно из трёх цифр;
    \item нормализация данных не менее важна, чем разбиение: без замены
    \enquote{ё} на \enquote{е} и удаления знака ударения одно слово становится
    двумя терминами или распадается на части;
    \item составные слова --- вопрос выбора: цельный токен
    (\enquote{санкт-петербург}) точнее, но словарь больше (978 тыс. против
    912 тыс. уникальных токенов при разбиении), поэтому разбиение вынесено в
    параметр;
    \item числа раздувают словарь: 53\% различных токенов встречаются ровно один раз (519 тыс.) --- в
    основном даты, десятичные дроби и номера; в индексе их, вероятно, стоит
    нормализовать или не индексировать;
    \item скорость регулярного выражения достаточная: 6{,}2 МБ/с (около
    500\,тыс.\ токенов в секунду) на одном потоке, весь корпус обработан за 35{,}5 с.
\end{itemize}

Известные ограничения: словарь сокращений конечен (незнакомое сокращение с точкой
остаётся словом, а точка отбрасывается); \enquote{г.} всегда считается сокращением,
даже в конце предложения; одиночная заглавная буква с точкой перед словом с
заглавной буквы считается инициалом; не обрабатываются \texttt{C++}, \texttt{C\#}
и другие символьные названия. Для задач поиска этого достаточно; полученный
токенизатор будет использоваться в следующих лабораторных работах при построении
индекса.
